% Draft. Compile with the IEEEtran class (IEEE Conference on Games); pdflatex paper.tex
% \todo{}: result not in yet.  \verify{}: number quoted from project notes, re-derive from the logs before submission.
\documentclass[conference]{IEEEtran}
\usepackage{amsmath,booktabs,xcolor,siunitx}
\newcommand{\todo}[1]{\textcolor{red}{[TODO: #1]}}
\newcommand{\verify}[1]{\textcolor{orange}{#1\textsuperscript{?}}}

\begin{document}
\title{Fitting the Leaf Evaluation of a Determinized-MCTS Agent\\for Pok\'emon Random Battles}
\author{\IEEEauthorblockN{Fernando Rimoli}}
\maketitle

\begin{abstract}
Determinized Monte Carlo tree search on a fast simulator is the strongest known approach to Pok\'emon Showdown
random battles; Foul Play, built on the poke-engine simulator, won the NeurIPS 2025 Pok\'emon battling challenge.
Its search scores leaves with a hand-written material evaluation. We fit that evaluation, 64 features in a
logistic model, to the outcomes of 40{,}000 engine self-play games, fitting within team pairs so that only the
variation a search compares identifies the weights. With nothing else changed, our search goes from
54.5\%~$\pm$~1.5\% to 59.2\%~$\pm$~2.2\% against Foul Play in Generation~4 random battles (95\% intervals,
1{,}600--3{,}200 games per arm), and still wins 57.8--58.4\% with compute matched in either direction. Rescaling
the stock evaluation alone gives no gain (54.5\%~$\pm$~2.3\%). \todo{Generation 9 replication.} With the same
statistical power we also report what did \emph{not} help: network priors, PUCT, regret matching, Bayesian root
solving, a value-network lookahead, more determinizations and longer self-play training.
\end{abstract}

\section{Introduction}
Pok\'emon battles are two-player, simultaneous-move, stochastic games with imperfect information: in random
battles each side's six Pok\'emon, their moves, items and abilities are drawn by a generator and hidden until
revealed. The strongest agents combine \emph{determinization} (sample the hidden information, search each sample as
a perfect-information game) with Monte Carlo tree search on a fast simulator~\cite{foulplay,pokeengine}.
Learned agents trained by reinforcement learning~\cite{wang2024} or offline from human replays~\cite{metamon}
are competitive in some formats but have not displaced search in random battles.

In such a search the leaf evaluation does most of the work: poke-engine does no rollouts, and every MCTS iteration
is scored by a hand-written function of material (alive Pok\'emon, HP, status, boosts, hazards, items) passed
through a sigmoid. We ask how much is gained by \emph{fitting} that function, as chess engines have long done
(``Texel tuning''), in a setting with simultaneous moves, chance and determinization.

\textbf{Contributions.}
(1) A fitted leaf evaluation that raises a determinized-MCTS agent from 54.5\% to 59.2\% against Foul Play at the
same compute, and an ablation showing the gain is in the weights, not the rescaling.
(2) A compute-matched comparison in both directions.
(3) A decomposition of where strength lies: the value of perfect information, the accuracy of set inference, and a
set of standard search and learning techniques that do not help, each measured with enough games to detect a
3--5 point effect.
(4) Benchmark methodology: why the 100--500-game evaluations common in this literature cannot resolve these
effects, and how little duplicate pairing helps.
(5) Open-source code, including fixes to the simulator's Generation~4 status durations.

\section{Background and related work}
\textbf{poke-engine and Foul Play.} poke-engine~\cite{pokeengine} implements battle mechanics, a decoupled-UCB1
simultaneous-move MCTS and a stock evaluation; MCTS scores a leaf $s$ as
$\sigma(0.0125\,(e(s)-e(s_{\text{root}})))$ with $e$ the material score. Foul Play~\cite{foulplay} samples opponent
sets from a set library, searches 2 determinizations for 1\,s each (4 for 0.5\,s early), sums visit counts and
plays among the moves within 75\% of the best. We use it only as an opponent.
\textbf{Evaluation tuning.} Logistic regression of game outcomes on evaluation features (Texel tuning), TD-Leaf and
TreeStrap~\cite{treestrap} fit evaluations in perfect-information games; we fit to outcomes of engine self-play.
\textbf{Learning agents.} Wang~\cite{wang2024} trains PPO by self-play for Generation~4 random battles (1693 Elo on
the public ladder) and adds MCTS at test time; Metamon~\cite{metamon} trains offline RL agents on millions of human
replays for OU formats.
\textbf{Imperfect information.} Determinization (perfect-information Monte Carlo) suffers strategy
fusion~\cite{frank1998}; information-set MCTS~\cite{ismcts} addresses it at a cost in sample efficiency.

\section{System}
\textbf{Search.} At each decision we draw $N{=}4$ determinizations, run poke-engine's MCTS on each for 1\,s in
parallel, average the root visit shares and play the argmax, optionally after adding $w\cdot\pi(a)$ for a policy
network $\pi$ (Section~\ref{sec:where}).
\textbf{Set inference.} Opponent sets are drawn from a library sampled from Showdown's own generator (2M teams),
in proportion to generator frequency, among the sets consistent with the battle so far: revealed moves, ability and
item; the lead slot; item evidence (no choice item after two different moves, no Leftovers after a turn damaged
without healing, no Life Orb after a hit without recoil); the speed order of each turn; and the damage of each hit
against the simulator's damage range for that set. Offline on true teams it identifies the whole set
\verify{59\%} of the time (item \verify{86--89\%}, moves \verify{79\%}, ability \verify{97\%}).
\textbf{Simulator corrections.} poke-engine~0.0.48 compiles out all Generation~4 taunt and encore durations and
never ends confusion; we branch on their random durations. This is a correctness fix and neutral for strength
(\verify{49.1\%~$\pm$~4.1\%} against the uncorrected engine).

\section{Fitting the leaf evaluation}\label{sec:fit}
\textbf{Features.} $f(s)\in\mathbb{R}^{64}$, each side one minus side two: $f_0$ is the stock score$/100$, the
others are corrections (alive counts, HP in quarter bins, active HP, each status and its HP interaction, items by
kind, Focus Sash intact, hazard damage per Pok\'emon, substitute, leech seed, confusion, taunt, yawn, locked moves,
recharge, each boost and its interaction with the Pok\'emon's attacking category, speed advantage, screens,
tailwind, wish, future sight, toxic counter).
\textbf{Data.} 39{,}823 games of engine self-play (Generation~4, 1{,}680{,}721 positions), both players searching
the full state with the stock evaluation for 20{,}000 iterations per move, a decision drawn from the visit shares
10\% of the time, every team pair played twice with sides swapped.
\textbf{Model.} $P(\text{win}\mid s)=\sigma(w^\top f(s))$, squared-probability loss, L2 on the corrections towards
the stock score (strength by cross-validation), splits by team pair. Because every leaf in one search shares the
teams, we also fit with a fixed effect per team pair, so that only within-pair variation identifies $w$.
The search uses $\sigma(k\,w^\top(f(s)-f(s_{\text{root}})))$.

\begin{table}[t]\centering\caption{Held-out log loss of the game outcome (lower is better).}\label{tab:fit}
\begin{tabular}{lcc}\toprule
Model & Gen.~4 & Gen.~9 \\\midrule
stock, engine scale & 0.526 & 0.536\\
stock, fitted scale & 0.522 & 0.527\\
all features & 0.511 & 0.514\\
all features, within team pair & 0.517 & 0.518\\\bottomrule
\end{tabular}\end{table}

\textbf{What the fit changes.} \todo{figure: fitted corrections by feature family, Gen.~4 vs Gen.~9.} The largest
Generation~4 corrections say the stock score overvalues stat boosts (attack, special attack, speed), undervalues
the cost of entry hazards and of being locked into a move, and undervalues a pending Wish; most corrections have
the same sign in Generation~9. \todo{interpret with feature scales.}
\textbf{Leaf scale.} The stock scale is too soft: with the fitted (within-pair) weights, doubling the scale
($k{=}2$) was chosen; at 300k iterations $k\in\{2,4,8\}$ score \verify{58/56/57\%} against the stock evaluation in
engine self-play, so the choice is not sensitive.

\section{Results}
\textbf{Protocol.} Local Showdown server, Generation~4 random battles, Foul Play at its default settings (one core
per instance), duplicate pairing (each pair of games uses the same seed with the teams swapped). We report 95\%
intervals from the variance of pair means; they are 5--10\% tighter than the binomial ones
(Section~\ref{sec:method}).

\begin{table}[t]\centering\caption{Win rate against Foul Play (Generation~4). Core-seconds per decision in
brackets.}\label{tab:fp}
\begin{tabular}{llrc}\toprule
Ours & Foul Play & Games & Win rate\\\midrule
stock eval, $4{\times}1$\,s [4] & default [2] & 3198 & 54.5\%~$\pm$~1.5\\
stock weights, scale $\times2$ [4] & default [2] & 1594 & 54.5\%~$\pm$~2.3\\
fitted eval [4] & default [2] & 1599 & \textbf{59.2\%~$\pm$~2.2}\\
fitted eval, $2{\times}1$\,s [2] & default [2] & 1597 & 57.8\%~$\pm$~2.2\\
fitted eval [4] & $4{\times}1$\,s [4] & 1605 & 58.4\%~$\pm$~2.2\\\bottomrule
\end{tabular}\end{table}

Table~\ref{tab:fp}: the fitted evaluation adds $+4.7\pm2.7$ points; rescaling the stock evaluation alone adds
nothing, so the gain is in the fitted weights. Halving our compute to Foul Play's costs about a point, and doubling
Foul Play's does not close the gap. In a direct comparison without Foul Play, the fitted search beats the stock
search \verify{54.9\%~$\pm$~1.6\%} on Showdown, so the gain is not specific to the opponent. Fitting to
luck-adjusted outcomes (a control-variate correction of chance) instead of raw outcomes did not help
(\verify{52.1\%~$\pm$~1.1\%} for raw in self-play), so the method is plain outcome regression.
All rows in Table~\ref{tab:fp} add a PPO policy prior with weight 0.25 to the visit shares (Section~\ref{sec:where}):
the 500M-step network in the first three rows, the 600M-step one in the compute-matched rows (the two do not differ
measurably). The Generation~9 search uses no network (prior weight 0): our network has no tera actions.

\textbf{Generation 9.} \todo{Foul Play benchmark, 3 arms $\times$ 1600 games, running.} The same recipe on 39{,}970
Generation~9 self-play games gives a larger held-out improvement (Table~\ref{tab:fit}); in engine self-play the
fitted evaluation at $k{=}2$ scores 60.2\%~$\pm$~1.4\% against the stock one and the rescaled stock evaluation
58.4\%~$\pm$~1.4\% (4000 games each).

\section{Where the strength is}\label{sec:where}
\textbf{Information.} Given the true sets of the opponent's revealed Pok\'emon, the same search scores
\verify{61.3\%} against Foul Play, given the whole true team \verify{65.4\%} (stock evaluation; baseline 54.5\%).
Most of the remaining gap is generator randomness that the battle cannot reveal. In Generation~9, where tera types
multiply the sets, inference identifies only \verify{43.7\%} of whole sets.

\textbf{Techniques that did not help} (stock evaluation, 400 games each unless noted, 95\% $\pm$4.9 points per
run; \verify{all}): 8 determinizations instead of 4 53.6\%; root PUCT on the network prior, $c{=}2$ 55.0\% (800
games), $c{=}8$ 53.2\%; network prior weight 0.5 54.2\%, 1.0 49.6\%; one-turn value-network lookahead 49.7\%;
regret matching in the tree 53.4\%; Bayesian root game, playing its most likely action 52.9\%, sampling from it
40.8\% (harmful: Foul Play does not exploit predictability, so mixing only plays worse moves); more PPO training
(500M and 600M steps against 400M) no gain; 40M further steps with boosted embedding learning ($\times30$,
$\times300$) 50.6\% and 48.5\% against the 600M network, unboosted 50.0\% (4000 raw-network games each,
$\pm$1.55).

\textbf{A methodological trap.} Labelling positions with a stronger search of the same engine (given the true
teams) cannot judge a learned component: such an oracle sided with the search in 82\% of the decisions where the
network and the search disagreed, yet in games a 0.25 network prior beat none by \verify{53.8\% to 46.2\%}.

\section{Benchmark methodology}\label{sec:method}
Single 400-game runs of one configuration ranged from 50.7\% to 57.8\%; a 3-point difference needs about 2000
games per arm. Duplicate pairing removes the variance due to the team draw, but in random battles most variance is
in play: it was worth only 1.1--1.4 times the games. Compute should be stated per decision in core-seconds;
Foul Play's default is 2 (two sequential 1\,s searches), ours 4 (four parallel ones).

\section{Limitations}
One strong opponent; no public rating (the official ladder flags bot accounts, the PokeAgent ladder has no random
battles). The self-play data were generated with the stock evaluation and fitted once; iterating, or targets from
deeper search (TreeStrap), may add more. Determinization's strategy fusion is untouched.

\section{Conclusion}
\todo{after the Generation 9 results.}

\bibliographystyle{IEEEtran}
\begin{thebibliography}{9}
\bibitem{foulplay} P.~Mariglia, ``Foul Play: a Pok\'emon Showdown battle bot,'' \url{https://github.com/pmariglia/foul-play}.
\bibitem{pokeengine} P.~Mariglia, ``poke-engine,'' \url{https://github.com/pmariglia/poke-engine}, v0.0.48.
\bibitem{wang2024} J.~Wang, ``Winning at Pok\'emon random battles using reinforcement learning,'' MEng thesis, MIT, 2024.
\bibitem{metamon} \todo{Metamon / PokeAgent paper, arXiv 2603.15563: authors, title.}
\bibitem{treestrap} J.~Veness, D.~Silver, A.~Blair, W.~Uther, ``Bootstrapping from game tree search,'' NeurIPS 2009.
\bibitem{frank1998} I.~Frank, D.~Basin, ``Search in games with incomplete information: a case study using Bridge card play,'' Artificial Intelligence, 1998.
\bibitem{ismcts} P.~Cowling, E.~Powley, D.~Whitehouse, ``Information set Monte Carlo tree search,'' IEEE TCIAIG, 2012.
\end{thebibliography}
\end{document}
